\documentclass[11pt,a4paper]{article}
\usepackage[left=23mm,right=23mm,top=22mm,bottom=22mm,footskip=11mm]{geometry}
\usepackage{amsmath,amssymb}
\usepackage[no-math]{fontspec}
\usepackage{xeCJK}
\setmainfont{Latin Modern Roman}
\setCJKmainfont{FandolSong-Regular.otf}[Path=assets/fonts/fandol/,BoldFont=FandolSong-Bold.otf,ItalicFont=FandolSong-Regular.otf]
\setCJKsansfont{FandolSong-Regular.otf}[Path=assets/fonts/fandol/,BoldFont=FandolSong-Bold.otf]
\xeCJKsetup{PunctStyle=quanjiao,CheckSingle=true}
\usepackage{setspace}\setstretch{1.10}
\setlength{\parindent}{2em}\setlength{\parskip}{4pt}
\usepackage{xcolor,titlesec,tcolorbox,tikz,graphicx,enumitem,longtable,booktabs,array,calc,needspace}
\tcbuselibrary{breakable,skins}\usetikzlibrary{arrows.meta,positioning}
\definecolor{navy}{HTML}{10136B}\definecolor{inkgray}{HTML}{474747}\definecolor{warm}{HTML}{94662C}\definecolor{teal}{HTML}{236B66}
\titleformat{\section}{\Large\bfseries\color{navy}}{\thesection}{0.7em}{}
\titleformat{\subsection}{\large\bfseries\color{navy}}{\thesubsection}{0.7em}{}
\titlespacing*{\section}{0pt}{17pt}{9pt}\titlespacing*{\subsection}{0pt}{12pt}{6pt}
\newtcolorbox{sboxenv}[2]{enhanced,breakable,colframe=#1,colback=#1!3!white,colbacktitle=#1,coltitle=white,fonttitle=\bfseries,title={#2},boxrule=0.45pt,arc=1.4pt,left=8pt,right=8pt,top=6pt,bottom=6pt,before skip=9pt,after skip=9pt,pad at break*=2mm}
\setlist[enumerate]{leftmargin=2.2em,itemsep=3pt,parsep=0pt,topsep=3pt}
\setlist[itemize]{label={},leftmargin=1em,itemsep=3pt}
\providecommand{\tightlist}{\setlength{\itemsep}{3pt}\setlength{\parskip}{0pt}}
\usepackage[unicode,hidelinks,linktoc=all]{hyperref}\usepackage{bookmark}
\hypersetup{pdftitle={二次曲线的矩阵理论与坐标变换},pdfauthor={Vincent Arqain}}
\renewcommand{\contentsname}{目录}\setcounter{tocdepth}{1}
\newcommand{\fig}[2]{\begin{center}\includegraphics[width=0.94\linewidth,height=0.33\textheight,keepaspectratio]{print/conics-matrix-figures/#1.pdf}\par\small #2\end{center}}
\widowpenalty=10000\clubpenalty=10000\displaywidowpenalty=10000\predisplaypenalty=5000
\allowdisplaybreaks[1]\raggedbottom\emergencystretch=2em
\setlength{\LTpre}{6pt}\setlength{\LTpost}{6pt}\renewcommand{\arraystretch}{1.18}
\begin{document}
\begin{center}{\LARGE\bfseries 二次曲线的矩阵理论与坐标变换\par}\vspace{7pt}
{\large 二次型矩阵、不变量、标准形与几何元素\par}\vspace{5pt}
{S 样式\qquad CM ams\qquad 扩写校订版}\end{center}
\textbf{主线：}将一般二次方程写成对称矩阵，用正交对角化确定坐标轴，用平移确定中心或顶点；在标准系中读出几何元素，再用同一个坐标变换写回原坐标系。保留原稿十节内容与三道例题，补全证明和适用条件，增加两道巩固例题。
\begin{center}
\begin{tikzpicture}[node distance=5mm and 4mm,box/.style={draw=navy!65,fill=navy!3,rounded corners=1pt,minimum width=32mm,minimum height=12mm,align=center,font=\small},>=Stealth]
\node[box](a){一般方程};\node[box,right=of a](b){矩阵与不变量};\node[box,right=of b](c){类型与特征方向};\node[box,right=of c](d){中心或顶点};
\node[box,below=of d](e){标准方程};\node[box,left=of e](f){半轴与离心率};\node[box,left=of f](g){焦点、轴、准线};\node[box,left=of g,fill=teal!5](h){原坐标表达};
\draw[->](a)--(b);\draw[->](b)--(c);\draw[->](c)--(d);\draw[->](d)--(e);\draw[->](e)--(f);\draw[->](f)--(g);\draw[->](g)--(h);
\end{tikzpicture}\end{center}
\textbf{记号：}$\boldsymbol x=(x,y)^T$，$X=(x,y,1)^T$；$Q=(\boldsymbol e_1\ \boldsymbol e_2)$，$\det Q=1$。常数项记为 $c_0$，焦半距记为 $c$；粗体 $\boldsymbol b$ 表示一次项系数列，普通 $b$ 表示半轴长。中文为 Fandol 宋体，数学采用 article 加载 amsmath、amssymb 的默认 Computer Modern 字体。
\begingroup\small\setlength{\parskip}{0pt}\tableofcontents\endgroup
\newpage
\Needspace{12\baselineskip}\section{二次曲线的矩阵表示}\label{ux4e8cux6b21ux66f2ux7ebfux7684ux77e9ux9635ux8868ux793a}

\Needspace{8\baselineskip}\subsection{一般方程与三种写法}\label{ux4e00ux822cux65b9ux7a0bux4e0eux4e09ux79cdux5199ux6cd5}

\begin{sboxenv}{navy}{核心公式：系数与矩阵}
设 \((a_{11},a_{12},a_{22})\ne(0,0,0)\)，一般二次方程为

\begin{equation*}
F(x,y)=a_{11}x^2+2a_{12}xy+a_{22}y^2+2a_{13}x+2a_{23}y+a_{33}=0.
\end{equation*}

定义

\begin{equation*}
A=\begin{pmatrix}a_{11}&a_{12}\\a_{12}&a_{22}\end{pmatrix},\quad \boldsymbol b=\begin{pmatrix}a_{13}\\a_{23}\end{pmatrix},\quad c_0=a_{33},
\end{equation*}

\begin{equation*}
\widehat A=\begin{pmatrix}A&\boldsymbol b\\\boldsymbol b^T&c_0\end{pmatrix}=\begin{pmatrix}a_{11}&a_{12}&a_{13}\\a_{12}&a_{22}&a_{23}\\a_{13}&a_{23}&a_{33}\end{pmatrix}.
\end{equation*}

于是

\begin{equation*}
F(\boldsymbol x)=\boldsymbol x^TA\boldsymbol x+2\boldsymbol b^T\boldsymbol x+c_0=X^T\widehat AX.\tag{1.1}
\end{equation*}

\end{sboxenv}

交叉项和一次项系数须先除以 \(2\)。例如 \(4xy\) 对应 \(a_{12}=2\)，而
\(-32x\) 对应 \(a_{13}=-16\)。

记 \(F_1=a_{11}x+a_{12}y+a_{13}\)，\(F_2=a_{12}x+a_{22}y+a_{23}\)，则

\begin{equation*}
\begin{pmatrix}F_1\\F_2\end{pmatrix}=A\boldsymbol x+\boldsymbol b=\frac12\nabla F,\qquad F(\boldsymbol x)=\boldsymbol x^T(A\boldsymbol x+\boldsymbol b)+\boldsymbol b^T\boldsymbol x+c_0.\tag{1.2}
\end{equation*}

\(\widehat A\) 决定所选方程；同一曲线也可由 \(\rho\widehat A\)
表示，其中
\(\rho\ne0\)。因此须区分``固定方程的坐标不变量''和``与方程整体倍乘无关的几何量''。

\Needspace{8\baselineskip}\subsection{不变量与特征方程}\label{ux4e0dux53d8ux91cfux4e0eux7279ux5f81ux65b9ux7a0b}

\begin{sboxenv}{navy}{核心公式：三个不变量与半不变量}
\begin{equation*}
\begin{aligned}I_1&=\operatorname{tr}A=a_{11}+a_{22},\\ I_2&=\det A=a_{11}a_{22}-a_{12}^2,\\ I_3&=\det\widehat A,\\ K_1&=\begin{vmatrix}a_{11}&a_{13}\\a_{13}&c_0\end{vmatrix}+\begin{vmatrix}a_{22}&a_{23}\\a_{23}&c_0\end{vmatrix}=I_1c_0-\boldsymbol b^T\boldsymbol b.\end{aligned}\tag{1.3}
\end{equation*}

特征方程为 \(\lambda^2-I_1\lambda+I_2=0\)。令

\begin{equation*}
\Delta=\sqrt{I_1^2-4I_2}=\sqrt{(a_{11}-a_{22})^2+4a_{12}^2},\qquad \lambda_\pm=\frac{I_1\pm\Delta}{2}.\tag{1.4}
\end{equation*}

\end{sboxenv}

两特征值恒为实数；\(\Delta=0\) 当且仅当 \(A\)
为非零数量矩阵。此时是圆型，但还需常数项的条件，才能得到实圆而非空集或一点。

若方程整体乘以 \(\rho\)，则 \(I_1,I_2,I_3,K_1\) 分别乘以
\(\rho,\rho^2,\rho^3,\rho^2\)。以下分类中使用的零性、\(I_2\) 与 \(K_1\)
的符号、\(I_1I_3\) 的符号均不受影响。

\Needspace{8\baselineskip}\subsection{几何对象的矩阵刻画}\label{ux51e0ux4f55ux5bf9ux8c61ux7684ux77e9ux9635ux523bux753b}

\begin{longtable}[]{@{}
  >{\raggedright\arraybackslash}p{(\columnwidth - 4\tabcolsep) * \real{0.3333}}
  >{\raggedright\arraybackslash}p{(\columnwidth - 4\tabcolsep) * \real{0.3333}}
  >{\raggedright\arraybackslash}p{(\columnwidth - 4\tabcolsep) * \real{0.3333}}@{}}
\toprule\noalign{}
\begin{minipage}[b]{\linewidth}\raggedright
对象
\end{minipage} & \begin{minipage}[b]{\linewidth}\raggedright
矩阵条件或方程
\end{minipage} & \begin{minipage}[b]{\linewidth}\raggedright
适用说明
\end{minipage} \\
\midrule\noalign{}
\endhead
\bottomrule\noalign{}
\endlastfoot
中心 \(\boldsymbol x_0\) & \(A\boldsymbol x_0+\boldsymbol b=0\) &
有解即给出方程的中心对称性 \\
渐近方向 \(\boldsymbol v\ne0\) & \(\boldsymbol v^TA\boldsymbol v=0\) &
是方向条件，不自动给出渐近线的位置 \\
主方向 & \(A\boldsymbol v=\lambda\boldsymbol v\) &
由实对称矩阵的特征方向给出 \\
共轭直径 & \(\boldsymbol v^T(A\boldsymbol x+\boldsymbol b)=0\) &
对非渐近弦方向，实弦中点落在这一直线上 \\
主轴 &
\(\lambda\boldsymbol v^T\boldsymbol x+\boldsymbol v^T\boldsymbol b=0\) &
\(\lambda\ne0\)，\(\boldsymbol v\) 是轴的法向量 \\
正则点切线 & \(X_0^T\widehat AX=0\) & \(F(\boldsymbol x_0)=0\) 且
\(\nabla F(\boldsymbol x_0)\ne0\) \\
中心形式 &
\(F=(\boldsymbol x-\boldsymbol x_0)^TA(\boldsymbol x-\boldsymbol x_0)+I_3/I_2\)
& \(I_2\ne0\) \\
\end{longtable}

\Needspace{5\baselineskip}\textbf{证明：} 将 \(\boldsymbol x=\boldsymbol x_0+\boldsymbol h\)
代入，得到

\begin{equation*}
F(\boldsymbol x_0+\boldsymbol h)-F(\boldsymbol x_0-\boldsymbol h)=4\boldsymbol h^T(A\boldsymbol x_0+\boldsymbol b).
\end{equation*}

右端对所有 \(\boldsymbol h\)
为零，恰好等价于中心条件。对非退化实曲线，这也刻画通常的几何中心。

沿方向 \(\boldsymbol v\) 的直线可写为
\(\boldsymbol m+t\boldsymbol v\)。展开得

\begin{equation*}
F(\boldsymbol m+t\boldsymbol v)=F(\boldsymbol m)+2t\boldsymbol v^T(A\boldsymbol m+\boldsymbol b)+t^2\boldsymbol v^TA\boldsymbol v.\tag{1.5}
\end{equation*}

若有两个不同交点，\(\boldsymbol m\)
是它们的中点当且仅当两根互为相反数，即一次项为零。由此得到共轭直径方程。反过来，直径上只有使二次方程有两个实根的点才对应实弦中点，例如椭圆的这部分轨迹一般只是直径中的一段。

若 \(A\boldsymbol v=\lambda\boldsymbol v\)，共轭直径的法向量就是
\(\boldsymbol v\)。在以 \(\boldsymbol v\)
为一条坐标轴的正交系中，平移消去这一坐标的一次项后，方程对该坐标变号不变，故这条直径是对称轴。切线公式则由

\begin{equation*}
\nabla F(\boldsymbol x_0)^T(\boldsymbol x-\boldsymbol x_0)=0
\end{equation*}

和 \(F(\boldsymbol x_0)=0\)
展开得到。奇异点不能直接套用这一公式作为唯一切线。

\Needspace{12\baselineskip}\section{坐标变换对方程与矩阵的作用}\label{ux5750ux6807ux53d8ux6362ux5bf9ux65b9ux7a0bux4e0eux77e9ux9635ux7684ux4f5cux7528}

\Needspace{8\baselineskip}\subsection{刚体变换与新轴方向}\label{ux521aux4f53ux53d8ux6362ux4e0eux65b0ux8f74ux65b9ux5411}

\begin{sboxenv}{navy}{核心公式：正变换与逆变换}
\begin{equation*}
Q=R(\theta)=\begin{pmatrix}\cos\theta&-\sin\theta\\\sin\theta&\cos\theta\end{pmatrix}=(\boldsymbol e_1\ \boldsymbol e_2),
\end{equation*}

\begin{equation*}
\boldsymbol x=Q\boldsymbol x'+\boldsymbol x_0,\qquad \boldsymbol x'=Q^T(\boldsymbol x-\boldsymbol x_0),\qquad T=\begin{pmatrix}Q&\boldsymbol x_0\\0^T&1\end{pmatrix},\quad X=TX'.\tag{2.1}
\end{equation*}

\end{sboxenv}

分量形式为

\begin{equation*}
\begin{aligned}x&=x'\cos\theta-y'\sin\theta+x_0,& x'&=(x-x_0)\cos\theta+(y-y_0)\sin\theta,\\ y&=x'\sin\theta+y'\cos\theta+y_0,& y'&=-(x-x_0)\sin\theta+(y-y_0)\cos\theta.\end{aligned}\tag{2.2}
\end{equation*}

\(Q\) 的列是新轴在旧坐标系中的方向，不能把列误读为行。\(\det Q=-1\)
时仍保长度和角度，但包含反射；调整一个特征向量的符号即可令
\(\det Q=1\)。

\Needspace{8\baselineskip}\subsection{变换律及不变量的证明}\label{ux53d8ux6362ux5f8bux53caux4e0dux53d8ux91cfux7684ux8bc1ux660e}

\begin{sboxenv}{navy}{定理 2.1：矩阵变换律}
在式（2.1）下，\(F(\boldsymbol x)=F'(\boldsymbol x')\)，且

\begin{equation*}
A'=Q^TAQ,\qquad \boldsymbol b'=Q^T(A\boldsymbol x_0+\boldsymbol b),\qquad c_0'=F(\boldsymbol x_0),\qquad \widehat A'=T^T\widehat AT.\tag{2.3}
\end{equation*}

\end{sboxenv}

\Needspace{5\baselineskip}\textbf{证明：} 直接展开：

\begin{equation*}
\begin{aligned}F(Q\boldsymbol x'+\boldsymbol x_0)&=\boldsymbol x'^TQ^TAQ\boldsymbol x'+2(A\boldsymbol x_0+\boldsymbol b)^TQ\boldsymbol x'+F(\boldsymbol x_0).
\end{aligned}
\end{equation*}

比较二次项、一次项和常数项即得前三式。最后一式也可由
\(X^T\widehat AX=X'^T(T^T\widehat AT)X'\) 得到。

\(A'=Q^{-1}AQ\)，所以 \(I_1,I_2\) 与特征值不变；\(\det T=1\)，所以
\(I_3\) 不变。允许反射时 \(\det T=-1\)，仍因 \((\det T)^2=1\) 而保持
\(I_3\)。

对纯旋转，\(K_1=I_1c_0-\|\boldsymbol b\|^2\) 不变。对平移
\(\boldsymbol x=\boldsymbol z+\boldsymbol h\)，记新量为 \(K_1'\)，由
\(A^2-I_1A+I_2I=0\) 得

\begin{equation*}
K_1'-K_1=I_2\|\boldsymbol h\|^2+2\boldsymbol h^T(I_1I-A)\boldsymbol b.\tag{2.4}
\end{equation*}

当 \(I_2=I_3=0\) 时，\(A\) 秩为 \(1\)。在特征坐标系中
\(A=\operatorname{diag}(0,I_1)\)，而 \(I_3=-I_1\beta_1^2=0\) 迫使
\(\beta_1=0\)，故
\(\boldsymbol b\in\operatorname{im}A\)，\((I_1I-A)\boldsymbol b=0\)。式（2.4）说明此时
\(K_1\) 在平移下也不变，这就是它在退化抛物型分类中的作用。

\Needspace{8\baselineskip}\subsection{旋转角的两种求法}\label{ux65cbux8f6cux89d2ux7684ux4e24ux79cdux6c42ux6cd5}

非对角元为

\begin{equation*}
a_{12}'=\frac12(a_{22}-a_{11})\sin2\theta+a_{12}\cos2\theta.
\end{equation*}

若 \(a_{12}\ne0\)，消去交叉项可用

\begin{equation*}
\cot2\theta=\frac{a_{11}-a_{22}}{2a_{12}}.\tag{2.5}
\end{equation*}

当 \(a_{11}=a_{22}\) 且 \(a_{12}\ne0\) 时，可取 \(\theta=\pi/4\)。若
\(a_{12}=0\)，原方程已无交叉项，不必使用含零分母的公式；若还满足
\(a_{11}=a_{22}\)，任何旋转都保持二次部分不变。

式（2.5）只确定两条轴，不决定哪一条作为 \(x'\) 轴。选定 \(x'\)
对应的特征值 \(\lambda_1\) 后，更稳妥的方法是求
\(A\boldsymbol e_1=\lambda_1\boldsymbol e_1\)；在 \(a_{12}\ne0\) 时有

\begin{equation*}
\tan\theta=\frac{\lambda_1-a_{11}}{a_{12}}.\tag{2.6}
\end{equation*}

再把 \(\boldsymbol e_1\) 逆时针旋转 \(90^\circ\) 得到
\(\boldsymbol e_2\)，便自动有 \(\det Q=1\)。

\Needspace{8\baselineskip}\subsection{主动与被动观点}\label{ux4e3bux52a8ux4e0eux88abux52a8ux89c2ux70b9}

\textbf{被动观点：} 曲线固定，把原点移到
\(\boldsymbol x_0\)，把坐标轴取为
\(\boldsymbol e_1,\boldsymbol e_2\)。同一个点的新坐标是
\(Q^T(\boldsymbol x-\boldsymbol x_0)\)。

\textbf{主动观点：} 坐标系固定，把标准位置的曲线 \(C_0\) 经
\(\varphi(\boldsymbol z)=Q\boldsymbol z+\boldsymbol x_0\) 搬到
\(C\)。于是
\(C=\varphi(C_0)\)，两条曲线全等，几何对象也随这一个刚体运动一起变换。

\Needspace{12\baselineskip}\section{化标准形与九类曲线}\label{ux5316ux6807ux51c6ux5f62ux4e0eux4e5dux7c7bux66f2ux7ebf}

\Needspace{8\baselineskip}\subsection{\texorpdfstring{有心型：\(I_2\ne0\)}{有心型：I\_2\textbackslash ne0}}\label{ux6709ux5fc3ux578bi_2ne0}

\begin{sboxenv}{navy}{定理 3.1：中心标准形}
令 \(\boldsymbol x_0=-A^{-1}\boldsymbol b\)，并取
\(Q^TAQ=\operatorname{diag}(\lambda_1,\lambda_2)\)。则

\begin{equation*}
F(\boldsymbol x_0)=\frac{I_3}{I_2},\qquad F(\boldsymbol x)=(\boldsymbol x-\boldsymbol x_0)^TA(\boldsymbol x-\boldsymbol x_0)+\frac{I_3}{I_2},
\end{equation*}

\begin{equation*}
\lambda_1x'^2+\lambda_2y'^2+\frac{I_3}{I_2}=0.\tag{3.1}
\end{equation*}

\end{sboxenv}

\Needspace{5\baselineskip}\textbf{证明：} 平移到中心使一次项为零。又由分块消元，

\begin{equation*}
\det\widehat A=\det A\,(c_0-\boldsymbol b^TA^{-1}\boldsymbol b)=I_2F(\boldsymbol x_0).
\end{equation*}

除以 \(I_2\)
即得常数项；再用实对称矩阵的正交对角化定理，得到式（3.1）。实际计算中解线性方程
\(A\boldsymbol x_0=-\boldsymbol b\) 通常比显式求逆更简洁。

\Needspace{8\baselineskip}\subsection{\texorpdfstring{无心型：\(I_2=0\)、\(I_3\ne0\)}{无心型：I\_2=0、I\_3\textbackslash ne0}}\label{ux65e0ux5fc3ux578bi_20i_3ne0}

\(A\ne0\) 且秩为 \(1\)，所以 \(I_1\ne0\)。取 \(\boldsymbol e_1\)
为零特征值的单位特征向量，\(\boldsymbol e_2\) 为 \(I_1\)
的单位特征向量，先只旋转，记
\(\boldsymbol z=Q^T\boldsymbol x\)，\(\beta_i=\boldsymbol e_i^T\boldsymbol b\)。方程变为

\begin{equation*}
I_1z_2^2+2\beta_1z_1+2\beta_2z_2+c_0=0,\qquad I_3=-I_1\beta_1^2.\tag{3.2}
\end{equation*}

于是 \(\beta_1\ne0\)，中心方程在第一分量要求 \(\beta_1=0\)，无解。配方得

\begin{equation*}
I_1\left(z_2+\frac{\beta_2}{I_1}\right)^2+2\beta_1\left(z_1-\frac{\beta_2^2/I_1-c_0}{2\beta_1}\right)=0.
\end{equation*}

\begin{sboxenv}{navy}{核心公式：顶点与抛物线标准形}
\begin{equation*}
z_{2,V}=-\frac{\beta_2}{I_1},\qquad z_{1,V}=\frac{\beta_2^2/I_1-c_0}{2\beta_1},\qquad \boldsymbol V=z_{1,V}\boldsymbol e_1+z_{2,V}\boldsymbol e_2.\tag{3.3}
\end{equation*}

选 \(\boldsymbol e_1\) 的符号使 \(I_1\beta_1<0\)，并令
\(x'=z_1-z_{1,V}\)、\(y'=z_2-z_{2,V}\)，则

\begin{equation*}
y'^2=2px',\qquad p=-\frac{\beta_1}{I_1}=\frac{|\beta_1|}{|I_1|}=\sqrt{-\frac{I_3}{I_1^3}}>0.\tag{3.4}
\end{equation*}

\end{sboxenv}

由于 \(\boldsymbol e_1^TA=0\)，任何平移都不改变
\(\boldsymbol e_1^T\boldsymbol b\)；因此开口方向的符号判断不依赖原点。

\Needspace{8\baselineskip}\subsection{\texorpdfstring{线心型：\(I_2=I_3=0\)}{线心型：I\_2=I\_3=0}}\label{ux7ebfux5fc3ux578bi_2i_30}

由式（3.2）有 \(\beta_1=0\)。沿 \(z_2\) 平移后，方程为

\begin{equation*}
I_1y'^2+\frac{K_1}{I_1}=0,\qquad y'^2=-\frac{K_1}{I_1^2}.\tag{3.5}
\end{equation*}

中心方程仅限制 \(z_2=-\beta_2/I_1\)，中心构成一条直线。\(K_1<0\)
时为两条不同的实平行直线；\(K_1=0\) 时为重合直线；\(K_1>0\)
时无实点，传统上称为一对虚平行直线。

\Needspace{8\baselineskip}\subsection{分类总表与圆的条件}\label{ux5206ux7c7bux603bux8868ux4e0eux5706ux7684ux6761ux4ef6}

表中 \(a,b,p>0\)，各标准形均在适当正交坐标系中书写。

\begin{longtable}[]{@{}
  >{\raggedright\arraybackslash}p{(\columnwidth - 6\tabcolsep) * \real{0.2500}}
  >{\raggedright\arraybackslash}p{(\columnwidth - 6\tabcolsep) * \real{0.2500}}
  >{\raggedright\arraybackslash}p{(\columnwidth - 6\tabcolsep) * \real{0.2500}}
  >{\raggedright\arraybackslash}p{(\columnwidth - 6\tabcolsep) * \real{0.2500}}@{}}
\toprule\noalign{}
\begin{minipage}[b]{\linewidth}\raggedright
二次部分
\end{minipage} & \begin{minipage}[b]{\linewidth}\raggedright
附加条件
\end{minipage} & \begin{minipage}[b]{\linewidth}\raggedright
曲线
\end{minipage} & \begin{minipage}[b]{\linewidth}\raggedright
标准形
\end{minipage} \\
\midrule\noalign{}
\endhead
\bottomrule\noalign{}
\endlastfoot
\(I_2>0\) & \(I_1I_3<0\) & 实椭圆 &
\(\dfrac{x'^2}{a^2}+\dfrac{y'^2}{b^2}=1\) \\
\(I_2>0\) & \(I_1I_3>0\) & 虚椭圆（无实点） &
\(\dfrac{x'^2}{a^2}+\dfrac{y'^2}{b^2}=-1\) \\
\(I_2>0\) & \(I_3=0\) & 一点 &
\(\dfrac{x'^2}{a^2}+\dfrac{y'^2}{b^2}=0\) \\
\(I_2<0\) & \(I_3\ne0\) & 双曲线 &
\(\dfrac{x'^2}{a^2}-\dfrac{y'^2}{b^2}=1\) \\
\(I_2<0\) & \(I_3=0\) & 两条相交直线 &
\(\dfrac{x'^2}{a^2}-\dfrac{y'^2}{b^2}=0\) \\
\(I_2=0\) & \(I_3\ne0\) & 抛物线 & \(y'^2=2px'\) \\
\(I_2=I_3=0\) & \(K_1<0\) & 两条平行直线 & \(y'^2=a^2\) \\
\(I_2=I_3=0\) & \(K_1=0\) & 重合直线 & \(y'^2=0\) \\
\(I_2=I_3=0\) & \(K_1>0\) & 虚平行直线（无实点） & \(y'^2=-a^2\) \\
\end{longtable}

\textbf{分类依据：} \(I_2>0\) 时两特征值同号，符号与 \(I_1\)
相同；将式（3.1）乘以该符号即可判断右端是正、负还是零。\(I_2<0\)
时两特征值异号，常数项非零给出双曲线，为零则二次式分解成两条相交直线。秩为
\(1\) 的情况已由式（3.4）、（3.5）穷尽。

实圆的充要条件是 \(I_1^2=4I_2\) 且
\(I_1I_3<0\)。``一点是一对共轭虚直线''是在复数域中分解二次式的说法，不能读成实平面中还有两条直线。

\begin{sboxenv}{warm}{辨析：配方法与正交变换}
一般配方法对应 \(\boldsymbol x=M\boldsymbol z+\boldsymbol h\)，得到
\(A'=M^TAM\)。\(M\) 可逆便保持惯性，但若 \(M\)
不正交，不能直接把新坐标中的长度当成原平面的长度。此时欧氏度量变为
\(d\boldsymbol z^TM^TM\,d\boldsymbol z\)。

因此，一般非正交配方适合判型；求焦点、轴长与离心率时应采用正交变换，或者同时正确追踪度量。纯平移配方当然保持长度，不能笼统说``配方法不能求几何元素''。

\end{sboxenv}

\Needspace{12\baselineskip}\section{椭圆与双曲线的几何元素}\label{ux692dux5706ux4e0eux53ccux66f2ux7ebfux7684ux51e0ux4f55ux5143ux7d20}

\Needspace{8\baselineskip}\subsection{焦轴特征值与半轴长}\label{ux7126ux8f74ux7279ux5f81ux503cux4e0eux534aux8f74ux957f}

焦点所在的轴称为焦轴。记它对应的特征值为 \(\lambda_1\)，方向为
\(\boldsymbol e_1\)；另一个特征值为 \(\lambda_2\)。椭圆的 \(\lambda_1\)
取绝对值较小者；双曲线的 \(\lambda_1\) 取与 \(I_3\) 同号者。

\begin{longtable}[]{@{}lll@{}}
\toprule\noalign{}
几何量 & 实椭圆 & 双曲线 \\
\midrule\noalign{}
\endhead
\bottomrule\noalign{}
\endlastfoot
\(a^2\) & \(-\dfrac{I_3}{I_2\lambda_1}\) &
\(-\dfrac{I_3}{I_2\lambda_1}\) \\
\(b^2\) & \(-\dfrac{I_3}{I_2\lambda_2}\) &
\(\dfrac{I_3}{I_2\lambda_2}\) \\
\(c^2\) & \(a^2-b^2\) & \(a^2+b^2\) \\
离心率 & \(e=c/a\) & \(e=c/a\) \\
\end{longtable}

\textbf{理由：}
在式（3.1）中将常数项移到右端。椭圆同号系数的绝对值越小，分母越大，故对应长轴。双曲线
\(I_2<0\)，右端 \(-I_3/I_2\) 与 \(I_3\)
同号，实轴系数须与右端同号。两种选择都保证 \(a^2,b^2>0\)。

\Needspace{8\baselineskip}\subsection{从标准坐标读出几何元素}\label{ux4eceux6807ux51c6ux5750ux6807ux8bfbux51faux51e0ux4f55ux5143ux7d20}

\begin{sboxenv}{navy}{核心公式：回到原坐标系的字典}
\begin{equation*}
x'=\boldsymbol e_1^T(\boldsymbol x-\boldsymbol x_0),\qquad y'=\boldsymbol e_2^T(\boldsymbol x-\boldsymbol x_0).\tag{4.1}
\end{equation*}

焦点为 \(\boldsymbol x_0\pm c\boldsymbol e_1\)；焦轴为
\(\boldsymbol e_2^T(\boldsymbol x-\boldsymbol x_0)=0\)；另一条主轴为
\(\boldsymbol e_1^T(\boldsymbol x-\boldsymbol x_0)=0\)。

焦轴顶点为 \(\boldsymbol x_0\pm a\boldsymbol e_1\)，椭圆另有短轴顶点
\(\boldsymbol x_0\pm b\boldsymbol e_2\)。非圆情形的准线为

\begin{equation*}
\boldsymbol e_1^T(\boldsymbol x-\boldsymbol x_0)=\pm\frac{a^2}{c}.\tag{4.2}
\end{equation*}

\end{sboxenv}

正号焦点与正号准线配对，负号同理。若
\(\boldsymbol e_1=(\cos\theta,\sin\theta)^T\)，焦点的分量就是
\((x_0\pm c\cos\theta,y_0\pm c\sin\theta)\)；焦轴为
\(-(x-x_0)\sin\theta+(y-y_0)\cos\theta=0\)。

双曲线的渐近线可直接由

\begin{equation*}
F(\boldsymbol x)-\frac{I_3}{I_2}=(\boldsymbol x-\boldsymbol x_0)^TA(\boldsymbol x-\boldsymbol x_0)=0\tag{4.3}
\end{equation*}

分解得到；在标准系中就是 \(x'/a\pm y'/b=0\)。圆的
\(c=0\)，两焦点重合于圆心，式（4.2）不适用；实圆没有有限的焦点---准线表示。

\Needspace{8\baselineskip}\subsection{不必写出完整变换的主轴公式}\label{ux4e0dux5fc5ux5199ux51faux5b8cux6574ux53d8ux6362ux7684ux4e3bux8f74ux516cux5f0f}

设 \(\boldsymbol v=(l,m)^T\) 是非零特征值 \(\lambda\) 的特征向量，则

\begin{equation*}
lF_1+mF_2=0\tag{4.4}
\end{equation*}

是法向量为 \(\boldsymbol v\) 的那条主轴，而非方向为 \(\boldsymbol v\)
的轴。这里不必单位化；当 \(a_{12}\ne0\) 时可取
\((l,m)=(a_{12},\lambda-a_{11})\)。它仍需要特征方向，但无需先求完整旋转矩阵、再代换全部坐标。

\Needspace{8\baselineskip}\subsection{离心率与椭圆面积}\label{ux79bbux5fc3ux7387ux4e0eux692dux5706ux9762ux79ef}

\begin{sboxenv}{navy}{定理 4.1：离心率的不变量表达}
对实椭圆或双曲线，按焦轴选择特征值后，

\begin{equation*}
e^2=1-\frac{\lambda_1}{\lambda_2}=\frac{2\Delta}{\Delta-\operatorname{sgn}(I_3)I_1}.\tag{4.5}
\end{equation*}

实椭圆面积为

\begin{equation*}
S=\pi ab=\frac{\pi|I_3|}{I_2^{3/2}}.\tag{4.6}
\end{equation*}

\end{sboxenv}

\Needspace{5\baselineskip}\textbf{证明：} 椭圆 \(b^2/a^2=\lambda_1/\lambda_2\)，故
\(e^2=1-b^2/a^2\)；双曲线 \(b^2/a^2=-\lambda_1/\lambda_2\)，故
\(e^2=1+b^2/a^2\)，都得到第一式。由焦轴选择规则可统一写成

\begin{equation*}
\lambda_1=\frac{I_1+\operatorname{sgn}(I_3)\Delta}{2},\qquad \lambda_2=\frac{I_1-\operatorname{sgn}(I_3)\Delta}{2},
\end{equation*}

代入即得第二式。对实椭圆，\(a^2b^2=I_3^2/(I_2^2\lambda_1\lambda_2)=I_3^2/I_2^3\)，于是得到面积公式。

圆对应 \(e=0\)；非圆实椭圆有 \(0<e<1\)；双曲线有
\(e>1\)。抛物线的轴向特征值为零，焦点---准线推导同样给出
\(e=1\)，但有心型的式（4.5）第二式不用于抛物线。

\Needspace{12\baselineskip}\section{抛物线的几何元素}\label{ux629bux7269ux7ebfux7684ux51e0ux4f55ux5143ux7d20}

\Needspace{8\baselineskip}\subsection{轴、顶点与开口方向}\label{ux8f74ux9876ux70b9ux4e0eux5f00ux53e3ux65b9ux5411}

设 \(I_2=0\)、\(I_3\ne0\)。取单位向量 \(\boldsymbol u\) 使
\(A=I_1\boldsymbol u\boldsymbol u^T\)；必须说明 \(\boldsymbol u\)
已单位化，不能只取任意比例的非零特征向量写此等式。

\(A\) 的任一非零列都平行于 \(\boldsymbol u\)。轴向
\(\boldsymbol e_1\perp\boldsymbol u\)，按
\(I_1\boldsymbol e_1^T\boldsymbol b<0\) 选择符号，则 \(\boldsymbol e_1\)
指向开口方向。

\begin{sboxenv}{navy}{核心公式：轴方程及其特殊情形}
轴为 \(\boldsymbol u^T(A\boldsymbol x+\boldsymbol b)=0\)。若
\((a_{11},a_{12})\ne(0,0)\)，可写成

\begin{equation*}
a_{11}F_1+a_{12}F_2=0.\tag{5.1}
\end{equation*}

若第一列为零，则必须改用第二列：

\begin{equation*}
a_{12}F_1+a_{22}F_2=0.\tag{5.2}
\end{equation*}

\end{sboxenv}

原稿的式（5.1）隐含第一列非零；例如 \(y^2-4x=0\)
的第一列为零，用它只能得到恒等式 \(0=0\)。第二列给出 \(y=0\)，这才是轴。

轴与曲线的唯一交点就是顶点
\(\boldsymbol V\)；也可直接用式（3.3）。在标准系中轴是 \(y'=0\)，代入
\(y'^2=2px'\) 得 \(x'=0\)，所以交点唯一。

\Needspace{8\baselineskip}\subsection{焦点、准线与顶点切线}\label{ux7126ux70b9ux51c6ux7ebfux4e0eux9876ux70b9ux5207ux7ebf}

\begin{sboxenv}{navy}{定理 5.1：抛物线几何元素}
令 \(p=\sqrt{-I_3/I_1^3}>0\)，则

\begin{equation*}
\boldsymbol F=\boldsymbol V+\frac p2\boldsymbol e_1,\qquad \text{准线：}\ \boldsymbol e_1^T(\boldsymbol x-\boldsymbol V)=-\frac p2,
\end{equation*}

\begin{equation*}
\text{顶点切线：}\ \boldsymbol e_1^T(\boldsymbol x-\boldsymbol V)=0.\tag{5.3}
\end{equation*}

\end{sboxenv}

\Needspace{5\baselineskip}\textbf{证明：} 在标准系 \(y'^2=2px'\) 中，点到 \((p/2,0)\) 的距离平方为

\begin{equation*}
\left(x'-\frac p2\right)^2+y'^2=\left(x'+\frac p2\right)^2,
\end{equation*}

恰等于到直线 \(x'=-p/2\)
的距离平方。焦点---准线定义给出焦点和准线，离心率为 \(1\)；在原点处对
\(y'^2-2px'\) 求梯度，切线为
\(x'=0\)。用式（4.1）的顶点版本拉回，即得结论。

这里采用 \(y'^2=2px'\) 的约定，顶点到焦点的距离为 \(p/2\)；若教材采用
\(y'^2=4fx'\)，则 \(f=p/2\)。

\Needspace{12\baselineskip}\section{原曲线与标准曲线的变换关系}\label{ux539fux66f2ux7ebfux4e0eux6807ux51c6ux66f2ux7ebfux7684ux53d8ux6362ux5173ux7cfb}

\Needspace{8\baselineskip}\subsection{标准系与原坐标系的对照}\label{ux6807ux51c6ux7cfbux4e0eux539fux5750ux6807ux7cfbux7684ux5bf9ux7167}

\begin{longtable}[]{@{}
  >{\raggedright\arraybackslash}p{(\columnwidth - 4\tabcolsep) * \real{0.3333}}
  >{\raggedright\arraybackslash}p{(\columnwidth - 4\tabcolsep) * \real{0.3333}}
  >{\raggedright\arraybackslash}p{(\columnwidth - 4\tabcolsep) * \real{0.3333}}@{}}
\toprule\noalign{}
\begin{minipage}[b]{\linewidth}\raggedright
对象
\end{minipage} & \begin{minipage}[b]{\linewidth}\raggedright
标准系
\end{minipage} & \begin{minipage}[b]{\linewidth}\raggedright
原坐标系
\end{minipage} \\
\midrule\noalign{}
\endhead
\bottomrule\noalign{}
\endlastfoot
点 & \(\boldsymbol P'\) &
\(\boldsymbol P=Q\boldsymbol P'+\boldsymbol x_0\) \\
方向 & \(\boldsymbol v'\) & \(\boldsymbol v=Q\boldsymbol v'\) \\
直线 & \(\boldsymbol n'^T\boldsymbol x'=d\) &
\((Q\boldsymbol n')^T(\boldsymbol x-\boldsymbol x_0)=d\) \\
曲线 & \(G(\boldsymbol x')=0\) &
\(G(Q^T(\boldsymbol x-\boldsymbol x_0))=0\) \\
齐次矩阵 & \(\widehat A'\) & \(\widehat A=T^{-T}\widehat A'T^{-1}\) \\
度量量 & 长度、角度、面积、\(a,b,c,e,p\) & 刚体运动下不变 \\
\end{longtable}

\textbf{法向量为何也乘 \(Q\)？} 将逆变换代入直线方程，有
\(\boldsymbol n'^TQ^T(\boldsymbol x-\boldsymbol x_0)=d\)。一般可逆仿射变换
\(\boldsymbol x=M\boldsymbol x'+\boldsymbol h\) 下，法向量按
\(M^{-T}\boldsymbol n'\) 变换；只有正交情形
\(M^{-T}=M\)，才与普通方向向量采用同一公式。

\Needspace{8\baselineskip}\subsection{参数方程的搬运}\label{ux53c2ux6570ux65b9ux7a0bux7684ux642cux8fd0}

\begin{sboxenv}{navy}{核心公式：原坐标中的参数方程}
\begin{equation*}
\begin{aligned}\text{椭圆：}\quad &\boldsymbol x(t)=\boldsymbol x_0+a\cos t\,\boldsymbol e_1+b\sin t\,\boldsymbol e_2,\quad 0\le t<2\pi;\\ \text{双曲线：}\quad &\boldsymbol x(t)=\boldsymbol x_0\pm a\cosh t\,\boldsymbol e_1+b\sinh t\,\boldsymbol e_2,\quad t\in\mathbb R;\\ \text{抛物线：}\quad &\boldsymbol x(t)=\boldsymbol V+\frac{t^2}{2p}\boldsymbol e_1+t\boldsymbol e_2,\quad t\in\mathbb R.\end{aligned}\tag{6.1}
\end{equation*}

\end{sboxenv}

这些是坐标列的简写。例如椭圆第一式就是
\(x=x_0+a e_{11}\cos t+b e_{21}\sin t\)、\(y=y_0+a e_{12}\cos t+b e_{22}\sin t\)。双曲线的两个符号分别参数化两支，\(\cosh^2t-\sinh^2t=1\)
保证方程成立。

\Needspace{8\baselineskip}\subsection{伸缩、旋转与平移的分解}\label{ux4f38ux7f29ux65cbux8f6cux4e0eux5e73ux79fbux7684ux5206ux89e3}

单位圆 \(u^2+v^2=1\) 经过伸缩、旋转、平移后得到椭圆：

\begin{equation*}
\boldsymbol x=\boldsymbol x_0+Q\operatorname{diag}(a,b)\begin{pmatrix}u\\v\end{pmatrix},
\end{equation*}

\begin{equation*}
\begin{pmatrix}Q\operatorname{diag}(a,b)&\boldsymbol x_0\\0^T&1\end{pmatrix}=\begin{pmatrix}I&\boldsymbol x_0\\0^T&1\end{pmatrix}\begin{pmatrix}Q&0\\0^T&1\end{pmatrix}\begin{pmatrix}\operatorname{diag}(a,b)&0\\0^T&1\end{pmatrix}.\tag{6.2}
\end{equation*}

矩阵从右向左作用。次序一般不可交换：先旋转、后平移
\(\boldsymbol x_0\)，等价于先平移
\(Q^T\boldsymbol x_0\)、后旋转。双曲线可从 \(u^2-v^2=1\) 作同样分解。

\fig{transport}{伸缩、旋转和平移将单位圆变成例 1 的椭圆；焦点只在后两步刚体运动中随曲线一起搬运。}

抛物线可由 \(v^2=u\) 经 \((x',y')=(u/(2p),v)\) 得到
\(y'^2=2px'\)，但这是非等比伸缩，本身不足以证明所有抛物线相似。要证明相似性，应取等比伸缩
\((x',y')=(2pu,2pv)\)，同样得到
\(y'^2=2px'\)。于是任意两条非退化实抛物线都可由相似变换互相得到。

\Needspace{8\baselineskip}\subsection{正交、仿射与射影三个层次}\label{ux6b63ux4ea4ux4effux5c04ux4e0eux5c04ux5f71ux4e09ux4e2aux5c42ux6b21}

\begin{longtable}[]{@{}
  >{\raggedright\arraybackslash}p{(\columnwidth - 4\tabcolsep) * \real{0.3333}}
  >{\raggedright\arraybackslash}p{(\columnwidth - 4\tabcolsep) * \real{0.3333}}
  >{\raggedright\arraybackslash}p{(\columnwidth - 4\tabcolsep) * \real{0.3333}}@{}}
\toprule\noalign{}
\begin{minipage}[b]{\linewidth}\raggedright
变换层次
\end{minipage} & \begin{minipage}[b]{\linewidth}\raggedright
矩阵关系
\end{minipage} & \begin{minipage}[b]{\linewidth}\raggedright
保持的性质与分类
\end{minipage} \\
\midrule\noalign{}
\endhead
\bottomrule\noalign{}
\endlastfoot
正交与平移 & \(A'=Q^TAQ\) & 保持距离、角度及焦点等度量信息；标准形含
\(a,b,p\) \\
一般仿射 & \(A'=M^TAM\) &
保持平行、中点、面积比、中心、共轭关系；非退化实曲线分椭圆、双曲线、抛物线三类 \\
射影 & \(\widehat A'=P^T\widehat AP\) &
保持点线关联与交比；有实点的非退化实射影二次曲线只有一类 \\
\end{longtable}

仿射分类的三种代表可取
\(x^2+y^2=1\)、\(x^2-y^2=1\)、\(y^2=x\)。射影变换允许改变无穷远直线，因而打破椭圆、双曲线、抛物线之间的仿射区别。对有实点且
\(I_3\ne0\) 的二次曲线，\(\widehat A\)
是非退化不定矩阵；由惯性定理并允许方程整体乘以
\(-1\)，均可化到同一射影标准形。无实点的``虚椭圆''不在这个结论内。

焦点与离心率依赖欧氏度量，不是一般仿射或射影不变量。相似变换仍保持离心率并相应搬运焦点，但会按比例改变长度。

\Needspace{12\baselineskip}\section{例题与巩固}\label{ux4f8bux9898ux4e0eux5de9ux56fa}

\Needspace{8\baselineskip}\subsection{例 1：椭圆}\label{ux4f8b-1ux692dux5706}

\begin{sboxenv}{inkgray}{例 1（原稿）}
将 \(5x^2+4xy+8y^2-32x-56y+80=0\)
化为标准形，求其焦点、主轴、准线、顶点、面积和参数方程。

\end{sboxenv}

\Needspace{5\baselineskip}\textbf{解：} 系数矩阵及不变量为

\begin{equation*}
A=\begin{pmatrix}5&2\\2&8\end{pmatrix},\quad \boldsymbol b=\begin{pmatrix}-16\\-28\end{pmatrix},\quad c_0=80,\qquad I_1=13,\ I_2=36,\ I_3=-1296.
\end{equation*}

\(I_2>0\) 且 \(I_1I_3<0\)，故是椭圆。特征方程
\(\lambda^2-13\lambda+36=0\) 给出 \(\lambda_1=4\)、\(\lambda_2=9\)；取

\begin{equation*}
\boldsymbol e_1=\frac1{\sqrt5}\begin{pmatrix}2\\-1\end{pmatrix},\quad \boldsymbol e_2=\frac1{\sqrt5}\begin{pmatrix}1\\2\end{pmatrix},\quad Q=\frac1{\sqrt5}\begin{pmatrix}2&1\\-1&2\end{pmatrix}.
\end{equation*}

\(\det Q=1\)，旋转角 \(\theta=-\arctan(1/2)\approx-26.57^\circ\)，且
\(\tan2\theta=-4/3\)，与式（2.5）相符。解中心方程

\begin{equation*}
5x_0+2y_0=16,\qquad 2x_0+8y_0=28
\end{equation*}

得
\(\boldsymbol x_0=(2,3)^T\)，\(F(\boldsymbol x_0)=-36=I_3/I_2\)。正、逆变换为

\begin{equation*}
\begin{aligned}x&=2+\frac{2x'+y'}{\sqrt5},&x'&=\frac{2(x-2)-(y-3)}{\sqrt5},\\ y&=3+\frac{-x'+2y'}{\sqrt5},&y'&=\frac{(x-2)+2(y-3)}{\sqrt5}.\end{aligned}
\end{equation*}

代入得到 \(4x'^2+9y'^2=36\)，即

\begin{equation*}
\frac{x'^2}{9}+\frac{y'^2}{4}=1.\tag{7.1}
\end{equation*}

所以
\(a=3\)、\(b=2\)、\(c=\sqrt5\)、\(e=\sqrt5/3\)。由不变量公式，\(\Delta=5\)，\(e^2=10/(5+13)=5/9\)。

\textbf{焦点：} \((2,3)\pm(2,-1)\)，即 \(F_+(4,2)\)、\(F_-(0,4)\)。

\textbf{长轴与短轴：} 长轴为 \(x+2y-8=0\)，短轴为 \(2x-y-1=0\)。也可用
\(F_1+2F_2=9x+18y-72\) 与 \(2F_1-F_2=8x-4y-4\) 验证。

\textbf{准线：} \(x'=\pm9/\sqrt5\) 给出 \(2x-y-10=0\) 与
\(2x-y+8=0\)，分别对应 \(F_+\)、\(F_-\)。

\textbf{顶点与面积：} 四顶点为

\begin{equation*}
\left(2\pm\frac6{\sqrt5},3\mp\frac3{\sqrt5}\right),\qquad \left(2\pm\frac2{\sqrt5},3\pm\frac4{\sqrt5}\right),
\end{equation*}

面积 \(S=\pi ab=6\pi=\pi|I_3|/I_2^{3/2}\)。

\textbf{参数方程：}

\begin{equation*}
x=2+\frac{6\cos t+2\sin t}{\sqrt5},\qquad y=3+\frac{-3\cos t+4\sin t}{\sqrt5},\qquad 0\le t<2\pi.
\end{equation*}

\fig{ellipse}{例 1 的椭圆及新坐标轴：中心为 (2,3)，焦点为 (4,2)、(0,4)，虚线为准线。}

\Needspace{8\baselineskip}\subsection{例 2：抛物线}\label{ux4f8b-2ux629bux7269ux7ebf}

\begin{sboxenv}{inkgray}{例 2（原稿）}
求 \(x^2-2xy+y^2-4x-4y+8=0\) 的标准形及几何元素。

\end{sboxenv}

\Needspace{5\baselineskip}\textbf{解：}

\begin{equation*}
A=\begin{pmatrix}1&-1\\-1&1\end{pmatrix},\quad\boldsymbol b=\begin{pmatrix}-2\\-2\end{pmatrix},\quad c_0=8,\qquad I_1=2,\ I_2=0,\ I_3=-16.
\end{equation*}

故是抛物线。取轴向 \(\boldsymbol e_1=(1,1)^T/\sqrt2\)，并取
\(\boldsymbol e_2=(-1,1)^T/\sqrt2\)。有
\(I_1\boldsymbol e_1^T\boldsymbol b=-4\sqrt2<0\)，所以开口朝 \((1,1)\)
方向，旋转角为 \(45^\circ\)。

由式（5.1），轴为 \(F_1-F_2=2(x-y)=0\)。令 \(x=y\)，原方程化为
\(-8x+8=0\)，顶点为 \(V(1,1)\)。又

\begin{equation*}
p=\sqrt{\frac{16}{8}}=\sqrt2,\qquad F=V+\frac p2\boldsymbol e_1=\left(\frac32,\frac32\right).
\end{equation*}

准线为 \(x+y-1=0\)，顶点切线为 \(x+y-2=0\)。取

\begin{equation*}
x'=\frac{x+y-2}{\sqrt2},\qquad y'=\frac{-x+y}{\sqrt2},
\end{equation*}

标准形是 \(y'^2=2\sqrt2\,x'\)。反代后 \((x-y)^2/2=2(x+y-2)\)，乘以 \(2\)
恰为原方程，完成检验。

\Needspace{8\baselineskip}\subsection{例 3：双曲线}\label{ux4f8b-3ux53ccux66f2ux7ebf}

\begin{sboxenv}{inkgray}{例 3（原稿）}
求 \(7x^2+24xy+10x-24y-161=0\) 的标准形、焦点、主轴、准线与渐近线。

\end{sboxenv}

\Needspace{5\baselineskip}\textbf{解：}

\begin{equation*}
A=\begin{pmatrix}7&12\\12&0\end{pmatrix},\quad\boldsymbol b=\begin{pmatrix}5\\-12\end{pmatrix},\quad c_0=-161,\qquad I_1=7,\ I_2=-144,\ I_3=20736.
\end{equation*}

\(I_2<0\)、\(I_3\ne0\)，是双曲线。特征值为 \(16,-9\)；焦轴特征值取与
\(I_3\) 同号的 \(16\)，故取

\begin{equation*}
\boldsymbol e_1=\frac15\begin{pmatrix}4\\3\end{pmatrix},\qquad\boldsymbol e_2=\frac15\begin{pmatrix}-3\\4\end{pmatrix},\qquad\theta=\arctan\frac34\approx36.87^\circ.
\end{equation*}

此时 \(\tan2\theta=24/7\)。解 \(7x_0+12y_0=-5\)、\(12x_0=12\)，得中心
\((1,-1)\)，\(I_3/I_2=-144\)。坐标变换为

\begin{equation*}
x'=\frac{4(x-1)+3(y+1)}5,\qquad y'=\frac{-3(x-1)+4(y+1)}5.
\end{equation*}

标准形 \(16x'^2-9y'^2=144\)，即

\begin{equation*}
\frac{x'^2}{9}-\frac{y'^2}{16}=1.\tag{7.2}
\end{equation*}

于是 \(a=3\)、\(b=4\)、\(c=5\)、\(e=5/3\)；不变量验算给出
\(e^2=2\times25/(25-7)=25/9\)。

\textbf{焦点与顶点：} 焦点为 \((1,-1)\pm(4,3)\)，即
\(F_+(5,2)\)、\(F_-(-3,-4)\)；实轴顶点为
\((17/5,4/5)\)、\((-7/5,-14/5)\)。

\textbf{两轴：} 实轴 \(3x-4y-7=0\)，虚轴 \(4x+3y-1=0\)。原稿的快速检验
\(-3F_1+4F_2=27x-36y-63=0\) 正好给出实轴。

\textbf{准线：} \(x'=\pm9/5\) 给出 \(4x+3y-10=0\) 与 \(4x+3y+8=0\)。

\textbf{渐近线：}

\begin{equation*}
F-\frac{I_3}{I_2}=7(x-1)^2+24(x-1)(y+1)=(x-1)(7x+24y+17).
\end{equation*}

因此渐近线为 \(x=1\) 与 \(7x+24y+17=0\)。

\fig{parabola-hyperbola}{左为例 2 的抛物线，右为例 3 的双曲线；绿色为焦轴，红点为焦点，虚线为准线，灰线为双曲线渐近线。}

\Needspace{8\baselineskip}\subsection{例
4：退化类型的参数分界}\label{ux4f8b-4ux9000ux5316ux7c7bux578bux7684ux53c2ux6570ux5206ux754c}

\begin{sboxenv}{inkgray}{例 4（补充）}
对参数 \(\tau\in\mathbb R\)，分类
\((x+y)^2-4(x+y)+\tau=0\)，并求其中心集合。

\end{sboxenv}

\Needspace{5\baselineskip}\textbf{解：}
\(A=\begin{pmatrix}1&1\\1&1\end{pmatrix}\)，\(\boldsymbol b=(-2,-2)^T\)，故
\(I_1=2\)、\(I_2=I_3=0\)、\(K_1=2\tau-8\)。配方为

\begin{equation*}
(x+y-2)^2=4-\tau.
\end{equation*}

当 \(\tau<4\) 时为两条平行直线 \(x+y=2\pm\sqrt{4-\tau}\)；\(\tau=4\)
时为重合直线 \(x+y=2\)；\(\tau>4\) 时无实点。中心方程均为
\(x+y=2\)。取正交坐标 \(y'=(x+y-2)/\sqrt2\)，标准形
\(y'^2=(4-\tau)/2\)，与式（3.5）一致。

空集在集合意义下关于每一点对称；这里的``中心集合''特指保持该二次多项式不变的代数中心条件，不能将两种说法混淆。

\Needspace{8\baselineskip}\subsection{例
5：第一列为零时的抛物线}\label{ux4f8b-5ux7b2cux4e00ux5217ux4e3aux96f6ux65f6ux7684ux629bux7269ux7ebf}

\begin{sboxenv}{inkgray}{例 5（补充）}
求 \(y^2-4x+6y+5=0\) 的轴、顶点、焦点与准线。

\end{sboxenv}

\Needspace{5\baselineskip}\textbf{解：}
\(A=\operatorname{diag}(0,1)\)，\(\boldsymbol b=(-2,3)^T\)，\(I_1=1\)、\(I_2=0\)、\(I_3=-4\)，所以
\(p=2\)。第一列为零，应使用第二列得到轴 \(F_2=y+3=0\)。

配方 \((y+3)^2=4(x+1)\) 表明 \(V=(-1,-3)\)，开口沿正 \(x\) 方向。焦点为
\((0,-3)\)，准线为 \(x=-2\)，顶点切线为
\(x=-1\)。本例同时检验了轴公式的分支和焦参数中的因子 \(2\)。

\Needspace{12\baselineskip}\section{焦点的秩一判别法}\label{ux7126ux70b9ux7684ux79e9ux4e00ux5224ux522bux6cd5}

\Needspace{8\baselineskip}\subsection{从焦点---准线定义到矩阵分解}\label{ux4eceux7126ux70b9ux51c6ux7ebfux5b9aux4e49ux5230ux77e9ux9635ux5206ux89e3}

设 \(\widehat A\) 表示一条非退化、非圆的实二次曲线，候选焦点为
\(P_0=(p_0,q_0)\)。点到 \(P_0\) 的距离平方对应矩阵

\begin{equation*}
P_{P_0}=\begin{pmatrix}1&0&-p_0\\0&1&-q_0\\-p_0&-q_0&p_0^2+q_0^2\end{pmatrix},\qquad X^TP_{P_0}X=(x-p_0)^2+(y-q_0)^2.
\end{equation*}

\begin{sboxenv}{navy}{定理 8.1：焦点的秩一判别}
\(P_0\) 是该曲线的焦点，当且仅当存在 \(\mu\ne0\) 与
\(\ell=(\ell_1,\ell_2,\ell_3)^T\)，使

\begin{equation*}
P_{P_0}-\mu\widehat A=\ell\ell^T,\qquad (\ell_1,\ell_2)\ne(0,0).\tag{8.1}
\end{equation*}

此时配对准线为 \(\ell_1x+\ell_2y+\ell_3=0\)，离心率满足
\(e^2=\ell_1^2+\ell_2^2\)。右端须为秩 \(1\) 的半正定矩阵，仅验证秩 \(1\)
不够。

\end{sboxenv}

\Needspace{5\baselineskip}\textbf{证明：} 若式（8.1）成立，令 \(N=\sqrt{\ell_1^2+\ell_2^2}>0\)。在
\(F=0\) 上，

\begin{equation*}
|PP_0|^2=(\ell_1x+\ell_2y+\ell_3)^2=N^2d(P,L)^2.
\end{equation*}

反过来，这一距离等式与 \(F=0\) 等价，因为两边相减正好是 \(\mu F\)。故
\(P_0,L\) 给出离心率 \(N\) 的焦点---准线表示。

若 \(P_0\) 是焦点，将准线写成
\(\boldsymbol n^T\boldsymbol x=d\)，\(\|\boldsymbol n\|=1\)，则距离关系给出二次多项式
\(G=|PP_0|^2-e^2(\boldsymbol n^T\boldsymbol x-d)^2\)。它与 \(F\)
有相同的非退化实二次曲线，故 \(G=\mu F\)。取 \(\ell=e(n_1,n_2,-d)^T\)
即得式（8.1）。

这里``相同曲线的两个二次方程成比例''可直接验证，无需额外假设：先作可逆仿射变换将曲线化到三种代表形。圆形代表
\(x^2+y^2=1\) 上代入 \((\cos t,\sin t)\)，比较
\(1,\cos t,\sin t,\cos2t,\sin2t\)
的系数，得到任意消失的二次式都是其倍数；双曲线代表 \(uv=1\) 上代入
\((t,1/t)\)，乘以 \(t^2\) 后比较多项式系数；抛物线代表 \(y^2=x\) 上代入
\((t^2,t)\)，比较各次幂系数。三种情形都得到所需结论。

\Needspace{8\baselineskip}\subsection{离心率公式为何出现特征值之比}\label{ux79bbux5fc3ux7387ux516cux5f0fux4e3aux4f55ux51faux73b0ux7279ux5f81ux503cux4e4bux6bd4}

取式（8.1）的左上二阶块，有

\begin{equation*}
I-\mu A=\begin{pmatrix}\ell_1\\\ell_2\end{pmatrix}\begin{pmatrix}\ell_1&\ell_2\end{pmatrix}.\tag{8.2}
\end{equation*}

右端是非零秩一半正定矩阵，像空间沿准线法向，即焦轴方向；垂直焦轴的方向是它的零空间。故
\(1-\mu\lambda_2=0\)，\(\mu=1/\lambda_2\)，而沿焦轴的特征值为

\begin{equation*}
e^2=1-\mu\lambda_1=1-\frac{\lambda_1}{\lambda_2}.
\end{equation*}

对于抛物线 \(\lambda_1=0\)，立即得到
\(e=1\)。圆的两个特征值相等，左上块若降秩就只能为零，不能得到有限准线所需的非零法向量；因此圆须单独处理，不能不加限制地套用定理
8.1。把准线说成``趋于无穷远''只是一种极限解释，不是实平面中的直线方程。

\Needspace{8\baselineskip}\subsection{用例 1
验证秩一分解}\label{ux7528ux4f8b-1-ux9a8cux8bc1ux79e9ux4e00ux5206ux89e3}

对焦点 \(F_+(4,2)\)，取 \(\mu=1/9\)。直接相减得到

\begin{equation*}
P_{F_+}-\frac19\widehat A=\frac19\begin{pmatrix}4&-2&-20\\-2&1&10\\-20&10&100\end{pmatrix}=\ell\ell^T,\qquad \ell=\frac13\begin{pmatrix}2\\-1\\-10\end{pmatrix}.
\end{equation*}

因此准线为
\(2x-y-10=0\)，\(e^2=(4+1)/9=5/9\)，与标准形法完全一致。秩一方法揭示焦点的矩阵结构，但实际求完整几何元素时，正交标准形通常更直观。

\Needspace{12\baselineskip}\section{常见易错点与检查清单}\label{ux5e38ux89c1ux6613ux9519ux70b9ux4e0eux68c0ux67e5ux6e05ux5355}

\begin{sboxenv}{warm}{检查清单}
\begin{enumerate}
\def\labelenumi{\arabic{enumi}.}
\tightlist
\item
  写矩阵前，先把交叉项和一次项系数除以 \(2\)；\(F_1,F_2\)
  是偏导数的一半。
\item
  \(Q\) 的列是新轴方向；正变换用 \(Q\)，逆变换用 \(Q^T\)。求出矩阵后检查
  \(Q^TQ=I\)、\(\det Q=1\)。
\item
  求几何长度时必须将特征向量单位化；只求轴方程时则可用任意非零比例。
\item
  中心型用
  \(A\boldsymbol x_0=-\boldsymbol b\)；抛物线无中心，应求顶点。线心型的中心不唯一。
\item
  椭圆长轴取 \(|\lambda|\) 较小者；双曲线实轴取与 \(I_3\)
  同号者。主轴公式 \(lF_1+mF_2=0\) 中 \((l,m)\) 是法向量。
\item
  抛物线开口由 \(I_1\boldsymbol e_1^T\boldsymbol b<0\)
  确定；第一列为零时须改用第二列求轴。
\item
  \(y'^2=2px'\) 的焦距是 \(p/2\)。圆不能使用含 \(a^2/c\)
  的有限准线公式。
\item
  一般非正交配方改变度量；若不追踪
  \(M^TM\)，不能把新坐标中的轴长、焦点照搬回去。
\item
  不变量针对固定方程的坐标变化；整体倍乘方程会改变它们的数值。离心率等几何量应与这种倍乘无关。
\item
  正则点才有本讲的唯一切线公式；共轭直径上的所有点未必都是真实弦的中点。
\item
  焦点秩一判别还要求半正定及非零准线法向量，圆另论。相似性须由等比伸缩证明。
\item
  最后反代标准方程，检查焦点与轴、准线的配对；图形用作核验，不能代替代数证明。
\end{enumerate}

\end{sboxenv}

\textbf{解题顺序：}
写矩阵与不变量，先判类型；选焦轴或开口方向，求中心或顶点；写正交变换并反代验证；在标准系读出几何量；用同一逆坐标字典给出原坐标答案。若只要求某个对象，可选用轴方程、中心形式或秩一判别，避免无关计算。

\Needspace{12\baselineskip}\section{知识体系与空间推广}\label{ux77e5ux8bc6ux4f53ux7cfbux4e0eux7a7aux95f4ux63a8ux5e7f}

\Needspace{8\baselineskip}\subsection{六条知识线索}\label{ux516dux6761ux77e5ux8bc6ux7ebfux7d22}

\begin{longtable}[]{@{}
  >{\raggedright\arraybackslash}p{(\columnwidth - 4\tabcolsep) * \real{0.3333}}
  >{\raggedright\arraybackslash}p{(\columnwidth - 4\tabcolsep) * \real{0.3333}}
  >{\raggedright\arraybackslash}p{(\columnwidth - 4\tabcolsep) * \real{0.3333}}@{}}
\toprule\noalign{}
\begin{minipage}[b]{\linewidth}\raggedright
线索
\end{minipage} & \begin{minipage}[b]{\linewidth}\raggedright
核心对象
\end{minipage} & \begin{minipage}[b]{\linewidth}\raggedright
需要能够独立重建的结论
\end{minipage} \\
\midrule\noalign{}
\endhead
\bottomrule\noalign{}
\endlastfoot
矩阵表示 & \(A,\boldsymbol b,c_0,\widehat A\) &
系数约定、梯度、中心、直径、切线 \\
坐标变换 & \(Q,T\) & 变换律、正逆变换、旋转角、主动与被动观点 \\
标准形分类 & \(I_1,I_2,I_3,K_1\) & 有心、无心、线心三型及九类表 \\
几何元素 & 特征方向与 \(a,b,c,e,p\) &
焦轴选择、焦点、顶点、轴、准线、渐近线及面积 \\
变换关系 & 参数方程与矩阵分解 &
法向量变换、伸缩旋转平移、三种分类层次 \\
进一步推广 & 秩一分解、空间对称矩阵 &
焦点判别及三维二次曲面的标准化思路 \\
\end{longtable}

\Needspace{8\baselineskip}\subsection{空间二次曲面的同一方法}\label{ux7a7aux95f4ux4e8cux6b21ux66f2ux9762ux7684ux540cux4e00ux65b9ux6cd5}

空间二次曲面写成

\begin{equation*}
F(\boldsymbol x)=\boldsymbol x^TA\boldsymbol x+2\boldsymbol b^T\boldsymbol x+c_0=0,\qquad A=A^T\in\mathbb R^{3\times3},
\end{equation*}

其齐次矩阵 \(\widehat A\) 为四阶。仍先用三阶正交矩阵将 \(A\)
对角化，再沿非零特征方向平移消去一次项。若 \(A\) 可逆，中心为
\(-A^{-1}\boldsymbol b\)，中心值为 \(\det\widehat A/\det A\)；若 \(A\)
奇异，则先检查一次项在 \(\ker A\)
上的分量，非零时产生抛物型，零时可能产生柱面或退化曲面。

沿用空间常见编号，可记

\begin{equation*}
I_1=\operatorname{tr}A,\quad I_2=\frac{(\operatorname{tr}A)^2-\operatorname{tr}(A^2)}2,\quad I_3=\det A,\quad I_4=\det\widehat A.
\end{equation*}

这里的 \(I_2,I_3\)
已是三维约定，不能与本讲二维定义直接混用。方法相同，分类需要额外考虑秩、惯性及零特征方向上的一次项；焦点概念也不能仅靠把二维公式增加一个坐标得到。

\textbf{编校说明：}
保留原稿十节内容及原三例；新增不变量平移公式的证明、抛物线顶点显式公式、退化参数例与零列例。修正或补明了直径的实弦范围、切线正则性、方程倍乘、圆的有限准线例外、抛物线轴的零列分支和相似性证明。所有句子与段落均不使用前后装饰点；数学中的必要运算符保留其含义。

\end{document}
